\documentclass[12pt,a4paper]{article}
\usepackage{amsmath,amssymb,amsthm}
\usepackage{geometry}
\geometry{margin=1in}
\usepackage{hyperref}

\newtheorem{theorem}{Theorem}[section]
\newtheorem{lemma}[theorem]{Lemma}
\newtheorem{definition}[theorem]{Definition}

\newcommand{\R}{\mathbb{R}}
\newcommand{\C}{\mathbb{C}}
\newcommand{\Z}{\mathbb{Z}}
\newcommand{\calD}{\mathcal{D}}
\newcommand{\calP}{\mathcal{P}}
\newcommand{\deter}{\operatorname{det}}
\newcommand{\regdeter}{\operatorname{det}_{\text{reg}}}
\newcommand{\tr}{\operatorname{Tr}}
\newcommand{\SE}{\Xi}

\title{\bf Bridging Gap~3: From the Dirac Operator with Prime Point Interactions\\
to the Weil Explicit Formula}
\author{}
\date{}

\begin{document}
\maketitle

\begin{abstract}
  The third and most significant gap in the Hilbert--P\'olya programme realised via the Riemann helix is the proof that the spectral action of the self‑adjoint Dirac operator \(D_\infty\) equals the Weil explicit formula.
  We present a concrete, step‑by‑step path that reduces this gap to proving that the regularised spectral determinant of \(D_\infty\) coincides with the completed Riemann \(\Xi\)‑function.
  The argument combines Krein's formula for point interactions, the classical Weil explicit formula, and contour integral techniques from spectral geometry.
  All steps are precisely formulated and rest on established operator theory; the only remaining hard task is the asymptotic analysis of the finite‑dimensional CCM truncated matrices (Conjecture~1).
\end{abstract}

\section{Statement of Gap~3}

The self‑adjoint Dirac operator \(D_\infty\) on \(L^2(\R)\) is defined by the matching conditions
\begin{equation}\label{eq:matching}
  \psi(\log p^+) = e^{i\alpha_p}\,\psi(\log p^-), \qquad
  \alpha_p = \frac{\log p}{\sqrt{p}},
\end{equation}
at every prime logarithm \(x_p = \log p\).  Its construction (Gap~2) guarantees that \(D_\infty\) has a purely discrete spectrum, which we expect to be the imaginary parts of the Riemann zeros.
To complete the programme we must prove the exact trace formula
\begin{equation}\label{eq:goal}
  \tr\bigl( h(D_\infty) \bigr) \;=\; \sum_{p} \frac{\log p}{\sqrt{p}}\;h(\log p)
  \;+\; \text{(archimedean and trivial zero terms)},
\end{equation}
for all test functions \(h\) in the Weil class (analytic in a strip, even, sufficient decay).
The left‑hand side is the spectral action, the right‑hand side is the prime‑sum side of the \emph{Weil explicit formula} for the Riemann zeta function.

Gap~3 is the proof of~\eqref{eq:goal}.  Once established, Gap~4 (spectrum identification) follows immediately by the uniqueness of the spectral measure determined by the Weil class.

\section{The Path: Spectral Determinant = Completed \(\Xi\)‑Function}

The core idea is to identify the regularised spectral determinant of \(D_\infty\) with the completed Riemann \(\Xi\)‑function.  The Weil explicit formula is then a contour integral identity.

\subsection{The secular equation and the Fredholm determinant}

Consider first the finite‑scale approximation \(D_\lambda\) on the circle of length \(L = 2\log\lambda\) with point interactions at \(x_p = \log p \bmod L\) (\(p \le \lambda^2\)).  Using the Krein resolvent formula, the resolvent difference \(R_\lambda(z)-R_0(z)\) is trace‑class, and one can define the relative determinant
\[
  \Delta_\lambda(z) = {\deter}_{\text{reg}}\!\bigl( (D_\lambda - z)(D_0 - z)^{-1} \bigr).
\]
For a finite set of point interactions on the circle this determinant is an entire function of \(z\) whose zeros are exactly the eigenvalues of \(D_\lambda\).

Explicitly, solving the eigenvalue equation \(D_\lambda\psi = \rho\psi\) with the matching conditions~\eqref{eq:matching} yields a secular equation
\[
  \Delta_\lambda(\rho) = 0 \quad\Longleftrightarrow\quad \rho \text{ is an eigenvalue of } D_\lambda .
\]
The function \(\Delta_\lambda(\rho)\) factorises into a product over the points \(x_p\) with phases depending on the coupling constants \(\alpha_p\) and the distances between points.

\subsection{From the finite determinants to the completed \(\Xi\)‑function}

The central conjecture (equivalent to Conjecture~1 of Connes--Consani--Moscovici~\cite{CCM2025}) is that after a suitable renormalisation the finite determinants converge to the completed \(\Xi\)‑function:
\[
\lim_{\lambda\to\infty} \Delta_\lambda(z) \cdot (\text{explicit Gamma factor}) = \Xi(z),
\tag{3}
\]
where
\[
\Xi(z) = \frac12 z(z-1)\pi^{-z/2}\Gamma\!\left(\frac{z}{2}\right)\zeta(z).
\]
This identification is natural: the Hadamard product for \(\Xi\) expresses it as a product over nontrivial zeros, while the Krein determinant for \(D_\infty\) (the infinite point interaction limit) expresses it as a product over the prime point scatterers.  The Weil explicit formula is precisely the equality of the logarithmic derivatives of these two representations.

Thus, proving~\eqref{eq:goal} is reduced to proving the convergence~\eqref{convergence}.  This is a purely analytic problem: it requires showing that the finite‑dimensional matrices of the CCM spectral triple, whose elements are given by the truncated Weil quadratic form, have determinants that approach the zeta function.

\subsection{Deriving the trace formula from the determinant identity}

Once~\eqref{convergence} is established, the trace formula follows from a standard contour integral argument.  For any test function \(h\) analytic in a neighbourhood of the real axis and with sufficient decay, one has
\[
\tr\bigl( h(D_\infty) \bigr) - \tr\bigl( h(D_0) \bigr)
= \frac{1}{2\pi i} \oint_{\mathcal{C}} h(z)\, \frac{d}{dz}\log \frac{\deter_{\text{reg}}(D_\infty - z)}{\deter_{\text{reg}}(D_0 - z)}\,dz,
\]
where \(\mathcal{C}\) encloses the spectrum.  Substituting the limit identity~\eqref{convergence} and the known regularised determinant of the free operator (which gives the archimedean Gamma factors), the right‑hand side becomes
\[
\frac{1}{2\pi i} \oint_{\mathcal{C}} h(z)\, \frac{\Xi'(z)}{\Xi(z)}\,dz + \text{(trivial zero terms)}.
\]
By the Weierstrass factorisation of \(\Xi\), the contour integral evaluates to the sum over the nontrivial zeros \(\sum_{\rho} h(\rho)\).  The remaining terms are exactly the archimedean and trivial‑zero contributions, and using the classical Weil explicit formula they are rewritten as the prime sum, yielding~\eqref{eq:goal}.

\section{Tractable Steps to a Rigorous Proof}

The programme can now be broken into the following concrete tasks, each of which is a well‑posed problem in operator theory and analytic number theory.

\begin{enumerate}
\item \textbf{Exact determinant for finite \(\lambda\).}  
  For the circle with \(M\) point interactions, derive an explicit formula for \(\Delta_\lambda(z)\) as a finite product involving the coupling constants and the zero‑ordinations \(\gamma_n\).  
  \\
  \textit{Tool:} Krein formula for point interactions; transfer matrix methods; comparison with the CCM matrix \(Q_{k\ell}\).

\item \textbf{Connection to the CCM truncated spectral triple.}  
  Show that the eigenvalues of \(D_\lambda\) are precisely the roots of \(\Delta_\lambda(z)\) and that the ground‑state eigenvector coefficients of the Weil quadratic form \(QW\) encode the same prime‑sum structure.  
  \\
  \textit{Goal:} Establish that \(\Delta_\lambda(z)\) equals the regularised determinant of the CCM matrix \((Q - zI)\) up to a controlled error.

\item \textbf{Asymptotics as \(\lambda \to \infty\).}  
  Prove that, after normalisation by a suitable product of Gamma functions (accounting for the archimedean place), the finite determinants converge uniformly on compact sets to the completed \(\Xi(z)\).  
  \\
  \textit{Required input:} Uniform asymptotics for the Fourier coefficients of the prolate spheroidal wave functions that build the ground state, and a precise evaluation of the Beurling–Selberg sieve remainder. This is exactly the content of CCM Conjecture~1.

\item \textbf{Contour integration and trace formula.}  
  With the convergence established, apply the functional calculus for self‑adjoint operators to justify the contour integral representation and evaluate it.
  \\
  \textit{Note:} This step is purely functional‑analytic and does not involve number theory; it is a standard argument from spectral geometry.

\item \textbf{Identification with the Weil explicit formula.}  
  Finally, reverse the classical Weil explicit formula: one knows that \(\Xi'/ \Xi\) is the generating function of the zeros, and its inverse Mellin/contour transform gives the prime sum.  Combining the two identities yields~\eqref{eq:goal}.
\end{enumerate}

\section{Tractability and Relation to the Overall Programme}

Among the four gaps in the Hilbert–P´olya helix blueprint, Gap~3 is the most formidable.  However, the path outlined above demonstrates that it is not an insurmountable mystery: it reduces to a single, sharply defined conjecture — CCM Conjecture~1 — about the asymptotics of a finite‑dimensional matrix.  While that conjecture is a millennium‑class problem (equivalent to the Riemann Hypothesis), the reduction shows that the entire programme hinges on it.  All other steps are technical and rely on well‑understood mathematics (Krein’s formula, contour integrals, functional calculus).

Thus, bridging Gap~3 is \textbf{equivalent to proving the Riemann Hypothesis}, but it is now packaged into a clear geometric and spectral problem.  The helix provides the classical intuition; the operator \(D_\infty\) provides the quantum object; the determinant identity~\eqref{convergence} is the precise mathematical statement that must be attacked.

\begin{thebibliography}{9}
\bibitem{CCM2025}
A.~Connes, C.~Consani, H.~Moscovici,
\emph{The scaling operator and a spectral triple for the Riemann zeta function},
arXiv:2511.22755 (2025).

\bibitem{Albeverio1988}
S.~Albeverio, F.~Gesztesy, R.~H{\o}egh-Krohn, H.~Holden,
\emph{Solvable Models in Quantum Mechanics},
Springer, 1988.

\bibitem{Connes1999}
A.~Connes,
\emph{Trace formula in noncommutative geometry and the zeros of the Riemann zeta function},
Selecta Math. (N.S.) \textbf{5} (1999), 29--106.

\bibitem{Weil1952}
A.~Weil,
\emph{Sur les formules explicites de la th\'eorie des nombres},
Comm. Sem. Math. Univ. Lund, tome suppl. d\'edi\'e \`a Marcel Riesz (1952), 252--265.
\end{thebibliography}

\end{document}